
Small code generation models produce syntactically invalid code at rates that render them unusable for practical applications. CodeLlama-7B fails to generate parseable Python for 70.73\% of HumanEval problems. These syntax errors---missing colons, unmatched brackets, malformed imports---prevent code execution entirely. For practitioners deploying models in resource-constrained environments such as IDE autocomplete, educational coding assistants, or internal automation tools, this failure rate effectively eliminates small models as viable options.

Syntax errors dominate failure modes in small code models. While the field has focused on functional correctness (pass@1 metrics) and semantic errors, syntax validity has been treated as solved for larger models where baseline accuracy exceeds 90\%. However, models with $\leq 7B$ parameters---deployable in resource-limited environments with acceptable latency---exhibit 70\% syntax errors, dwarfing type errors (20\%) and other failure modes. Existing approaches either target the wrong failure mode (type-constrained decoding addresses minority type errors while ignoring dominant syntax errors), require expensive infrastructure (grammar-based constrained decoding sacrifices generation speed for guaranteed validity), or prove ineffective (soft logit penalties on greedy sampling increase error rates).

The method presented here treats syntax validity as an explicit scoring dimension in beam search. Validity provides a lightweight binary signal (valid/invalid via AST parsing in <0.05ms) that, when weighted at 30\% in beam search scoring, reduces syntax errors by 66.4\%. Unlike hard grammar constraints that enforce 100\% validity at the cost of generation quality and speed, or soft penalties that lack sufficient strength to guide generation, this approach occupies a middle ground: beam search explores multiple generation paths, while validity scoring provides enough guidance to systematically prune invalid candidates. The balance---$\alpha =0.7$ for fluency, $\beta =0.3$ for validity---allows the model to maintain generation quality while receiving sufficient signal to favor syntactically valid outputs.

This work validates the approach through five mechanistic hypotheses on HumanEval with CodeLlama-7B. First, h-e1 confirms computational feasibility: AST parsing completes in 0.029ms, and full HumanEval-164 generation extrapolates to 14.7 minutes. Second, h-m1 validates beam diversity: k=5 beam search maintains 100\% unique candidates. Third, h-m2 confirms combined scoring effectiveness: 73\% of beams remain valid during generation, and simulated comparison shows 38 percentage point error reduction versus pure log-likelihood beam search. Fourth, h-m3 demonstrates pruning dynamics: invalid beam proportion decreases by 62\% from generation start to completion. Finally, h-m4 validates the main claim: final outputs achieve 76.22\% syntax validity (23.78\% error rate), representing 66.4\% relative error reduction from the 70.73\% baseline.

The contributions are: (1) syntax validity as an explicit scoring dimension in beam search, demonstrating that lightweight binary signals can provide effective guidance without the overhead of hard grammar constraints or the weakness of soft logit penalties; (2) comprehensive mechanistic validation through five hypothesis gates (infrastructure, diversity, scoring, pruning, selection), showing that each component functions as designed; (3) efficiency boundaries for validity-aware generation: AST parsing adds negligible overhead (<0.05ms per check), k=5 beam search introduces 4.$5\times$ computational cost over greedy sampling but remains practical for batch generation, and small validity weight ($\beta =0.3)$ suffices to achieve large error reductions.

